\section{Control}
\label{sec:control}

\subsection{Why a cascade}

The vehicle is underactuated: four independent rotor speeds control six
degrees of freedom. Horizontal motion is available only by tilting the thrust
vector, so the causal chain is
\begin{equation}
  \tau_y \;\to\; \ddt\theta \;\to\; \theta \;\to\; \ddt x \;\to\; x ,
  \label{eq:chain}
\end{equation}
four integrations from moment to position. Near hover, linearising
\cref{eq:vdot,eq:wdot} about $R=I$, $T=mg$ gives the classical decoupling
\begin{equation}
  \ddt{x}\approx g\theta,\quad
  \ddt{y}\approx -g\phi,\quad
  \ddt{z}\approx \frac{\delta T}{m},\quad
  \ddt\phi = \frac{\tau_x}{J_x},\quad
  \ddt\theta = \frac{\tau_y}{J_y},\quad
  \ddt\psi = \frac{\tau_z}{J_z} .
  \label{eq:linhover}
\end{equation}
The attitude loop is two integrations from its input, the position loop four.
Since $\sqrt{K_R}\gg\sqrt{K_p}$ in any sensible design, time-scale separation
justifies the cascade

\begin{center}
position $\to$ velocity $\to$ attitude $\to$ body rate $\to$ rotor speed,
\end{center}

each level treating the level below as instantaneous. This is standard
practice~\cite{mahony2012}, and \cref{eq:linhover} is its justification.
\Cref{fig:cascade} shows the cascade as the engine implements it, and
\cref{fig:bandwidth} the separation of its time scales that the argument
needs.

\begin{figure}[tbp]
\centering
% The cascade of the engine's controller, from the user to the rigid body.
\begin{tikzpicture}[
    font=\footnotesize,
    box/.style={block, text width=1.95cm, minimum height=3.2em, inner sep=2pt},
    ctl/.style={box, fill=cA!7},
    phys/.style={box, fill=cB!10},
    sig/.style={font=\scriptsize, midway},
  ]
  % top row: from the user to the thrust command
  \node[phys] (user) at (0,0)      {user\\ $\vect c(t),\ \psi_h(t)$};
  \node[box]  (raw)  at (2.95,0)   {ideal display\\ trajectory};
  \node[ctl]  (gov)  at (5.9,0)    {reference\\ governor};
  \node[ctl]  (pos)  at (8.85,0)   {position loop\\ \labelcref{eq:poscmd}};
  \node[ctl]  (des)  at (11.8,0)   {desired attitude\\ \labelcref{eq:Rd,eq:omegad}};
  % bottom row: from the attitude loop back to the plant
  \node[ctl]  (att)  at (11.8,-2.6) {attitude loop\\ \labelcref{eq:leeused}};
  \node[ctl]  (mix)  at (8.85,-2.6) {mixer $\mat M^{+}$,\\ saturation};
  \node[phys] (mot)  at (5.9,-2.6)  {motors, lag\\ \labelcref{eq:motorlag}};
  \node[phys] (body) at (2.95,-2.6) {rigid body,\\ rotors, drag};
  \node[ctl]  (gim)  at (0,-2.6)    {gimbal PD,\\ screen upright};

  \draw[flow] (user) -- node[sig, above] {$\vect r_h$} (raw);
  \draw[flow] (raw)  -- node[sig, above] {$\vect r_{\mathrm{raw}}$} (gov);
  \draw[flow] (gov)  -- node[sig, above, align=center] {$\vect r_g,\vect v_g$\\ $\vect a_g,\vect j_g$} (pos);
  \draw[flow] (pos)  -- node[sig, above] {$\vect F_c$} (des);
  \draw[flow] (des)  -- node[sig, right, align=left] {$R_d$\\ $\vect\omega_d$} (att);
  \draw[flow] (att)  -- node[sig, above] {$\vect\tau$} (mix);
  \draw[flow] (mix)  -- node[sig, above] {$\Omega^{\mathrm{cmd}}$} (mot);
  \draw[flow] (mot)  -- node[sig, above] {$\Omega$} (body);
  % thrust magnitude goes to the mixer directly
  \draw[flow] ($(des.south)+(-0.7,0)$) -- ++(0,-0.62)
        -| node[sig, pos=0.25, above] {$T_c=\|\vect F_c\|$} ($(mix.north)+(0.55,0)$);
  % head position enters the keep-out constraint
  \draw[flow, cD] (user.north) -- ++(0,0.55)
        -| node[sig, pos=0.25, above] {head $\vect c,\dt{\vect c},\ddt{\vect c}$: keep-out \labelcref{eq:barrier}} (gov.north);
  % heading reference
  \draw[flow] ($(gov.north)+(0.65,0)$) -- ++(0,0.28)
        -| node[sig, pos=0.25, above] {$\psi_g,\ \dt\psi_g,\ \ddt\psi_g$} (des.north);
  % feedback of translational state
  \draw[flow, cE] ($(body.north)+(0.4,0)$) -- ++(0,0.5)
        -| node[sig, pos=0.2, below] {$\vect r,\ \vect v$} ($(pos.south)+(-0.45,0)$);
  % feedback of rotational state
  \draw[flow, cE] (body.south) -- ++(0,-0.8)
        -| node[sig, pos=0.25, below] {$R,\ \vect\omega$} (att.south);
  \draw[flow, cE] (body.west) -- node[sig, above] {$R$} (gim.east);
  \draw[flow] (gim.south) -- ++(0,-0.5)
        -| node[sig, pos=0.3, below] {$\vect\tau_g$} ($(body.south)+(-0.65,0)$);

  % the sampled controller
  \begin{scope}[on background layer]
    \node[draw=cA, dashed, rounded corners=3pt, inner xsep=6pt, inner ysep=9pt,
          fit=(gov)(pos)(des)(att)(mix)] (ctlbox) {};
  \end{scope}
  \node[font=\scriptsize, anchor=north west, align=left] at (-1.05,-4.3)
       {\tikz\fill[cA!7,draw=black] (0,0) rectangle (0.3,0.18); computed\quad
        \tikz\fill[cB!10,draw=black] (0,0) rectangle (0.3,0.18); physical\quad
        \tikz\draw[cA, dashed] (0,0) rectangle (0.3,0.18); controller, sampled at
        $1/h=\SI{250}{\hertz}$, output held between samples};
\end{tikzpicture}

\caption{The cascade as the engine implements it. Blue boxes are computed at
each sample of the controller, orange boxes are physical. The reference
governor turns the ideal display trajectory into one the aircraft can fly and
holds it outside the keep-out sphere about the user's head; the position loop
turns the governed reference into a force $\vect F_c$; the thrust magnitude,
the desired attitude $R_d$ and its rate $\vect\omega_d$ follow from that force
and its derivative; the attitude loop and the mixer turn them into rotor speed
commands. The translational state closes the outer loop, the rotational state
the inner one. The gimbal loop runs beside the cascade and keeps the screen
upright while the airframe tilts.}
\label{fig:cascade}
\end{figure}

\begin{remark}[Flatness]
The quadrotor with these dynamics is differentially flat with flat outputs
$(x,y,z,\psi)$~\cite{mellinger2011}: every state and input is an algebraic
function of the flat outputs and finitely many derivatives. This is what makes
feedforward from a desired trajectory possible and is why the controller below
uses $\ddt{\vect r}_d$ directly rather than relying on feedback alone.
\end{remark}

\subsection{Position loop}

Given a reference $(\vect r_d,\dt{\vect r}_d,\ddt{\vect r}_d)$, define errors
$\vect e_p = \vect r_d - \vect r$, $\vect e_v = \dt{\vect r}_d - \vect v$, and
command
\begin{equation}
  \vect a_c = \ddt{\vect r}_d + K_p\vect e_p + K_d\vect e_v
              + K_i\!\int\!\vect e_p\,\dd t ,
  \qquad
  \vect F_c = m(\vect a_c + g\ez) .
  \label{eq:poscmd}
\end{equation}
The desired thrust magnitude and thrust direction are
\begin{equation}
  T_c = \|\vect F_c\|,
  \qquad
  \vect b_3^d = \frac{\vect F_c}{\|\vect F_c\|} ,
  \label{eq:thrustdir}
\end{equation}
and with a desired yaw $\psi_d$ the desired attitude is completed by
\begin{equation}
  \vect b_2^d = \frac{\vect b_3^d \times \vect b_1^{\psi}}
                     {\|\vect b_3^d \times \vect b_1^{\psi}\|},
  \quad
  \vect b_1^d = \vect b_2^d\times\vect b_3^d,
  \quad
  R_d = [\,\vect b_1^d\ \vect b_2^d\ \vect b_3^d\,],
  \label{eq:Rd}
\end{equation}
where $\vect b_1^{\psi}=(\cos\psi_d,\sin\psi_d,0)$. \Cref{eq:Rd} is singular
only when $\vect b_3^d\parallel\vect b_1^{\psi}$, i.e.\ at \SI{90}{\degree} of
tilt, which the thrust limit precludes; the engine falls back to an alternative
reference axis if it occurs.

\subsection{Geometric attitude control on \texorpdfstring{$\SOthree$}{SO(3)}}

Attitude error is measured on the group rather than in coordinates, which
avoids both the singularities of Euler angles and the unwinding of naive
quaternion feedback.

\begin{definition}[Attitude error]
\begin{equation}
  \vect e_R = \tfrac12\big(R_d\transp R - R\transp R_d\big)^{\vee},
  \qquad
  \vect e_\omega = \vect\omega - R\transp R_d\,\vect\omega_d ,
  \label{eq:attitudeerror}
\end{equation}
with configuration error function
$\Psi(R,R_d) = \tfrac12\tr\!\big(I - R_d\transp R\big) \in [0,2]$.
\end{definition}

\begin{model}[Geometric attitude law]
\label{mod:leecontrol}
\begin{equation}
  \vect\tau = -K_R\,\vect e_R - K_\omega\,\vect e_\omega
              + \vect\omega\times J\vect\omega
              - J\big(\hatmap{\vect\omega}R\transp R_d\vect\omega_d
                      - R\transp R_d\dt{\vect\omega}_d\big) .
  \label{eq:leecontrol}
\end{equation}
\end{model}

\begin{theorem}[Exponential attitude stability; Lee, Leok and
McClamroch~\cite{lee2010}, Proposition~1]
\label{thm:lee}
Consider \cref{eq:wdot} under \cref{eq:leecontrol} with positive scalar gains
$k_R,k_\omega$ in place of $K_R,K_\omega$ and a smooth reference
$(R_d,\vect\omega_d,\dt{\vect\omega}_d)$. The equilibrium
$(\vect e_R,\vect e_\omega)=(0,0)$ is exponentially stable, and its region of
attraction contains the set
\begin{equation}
  \Psi(R(0),R_d(0)) < 2,
  \qquad
  \|\vect e_\omega(0)\|^2 < \frac{2}{\lambda_{\max}(J)}\,k_R\big(2-\Psi(R(0),R_d(0))\big) .
  \label{eq:roa}
\end{equation}
\end{theorem}

The excluded set is $\Psi=2$, where $R_d\transp R$ is a rotation by $\pi$. For
the unweighted $\Psi$ used here every such rotation is a critical point: if
$R_d\transp R = 2\vect u\vect u\transp - I$ for a unit vector $\vect u$, the
matrix is symmetric and $\vect e_R=0$. The critical set of $\Psi$ is therefore
the identity together with a two-dimensional family of attitudes, not three
isolated points. \Cref{thm:lee} is a statement on the sublevel set
\cref{eq:roa}, and no almost-global claim is made here. \Cref{fig:psi} shows
the mechanism along any single axis of rotation.

\begin{figure}[tbp]
\centering
% Configuration error Psi and the attitude error e_R against the rotation
% angle theta between R and R_d: Psi = 1 - cos theta, |e_R| = sin theta.
\begin{tikzpicture}
  \begin{axis}[paperplot, height=5.2cm,
      xmin=0, xmax=3.1416, ymin=0, ymax=2.25,
      xtick={0, 0.7854, 1.5708, 2.3562, 3.1416},
      xticklabels={$0$, $\pi/4$, $\pi/2$, $3\pi/4$, $\pi$},
      xlabel={rotation angle $\theta$ of $R_d\transp R$ about any axis $\vect u$},
      ylabel={dimensionless},
      legend pos=north west, domain=0:3.1416, samples=120]
    \addplot[name path=psi, cA] {1 - cos(deg(x))};
    \addlegendentry{$\Psi = \tfrac12\tr(I-R_d\transp R) = 1-\cos\theta$}
    \addplot[name path=er, cB] {sin(deg(x))};
    \addlegendentry{$\|\vect e_R\| = \sin\theta$}
    \path[name path=axisbot] (axis cs:0,0) -- (axis cs:3.1416,0);
    \addplot[cB!12] fill between[of=er and axisbot, soft clip={domain=1.5708:3.1416}];
    \draw[construct] (axis cs:1.5708,0) -- (axis cs:1.5708,2.25);
    \draw[construct, cD] (axis cs:0,2) -- (axis cs:3.1416,2);
    \node[lbl, cD, anchor=south east, align=right] at (axis cs:3.05,2.0)
      {$\Psi=2$: excluded by \cref{eq:roa}};
    \fill[cD] (axis cs:3.1416,0) circle (2.2pt);
    \fill[cD] (axis cs:3.1416,2) circle (2.2pt);
    \node[lbl, align=center, anchor=south] at (axis cs:2.2,0.08)
      {restoring term $K_R\vect e_R$\\ weakens past $\pi/2$};
    \draw[lbl, cD, -{Stealth[length=1.6mm]}] (axis cs:2.8,0.95) -- (axis cs:3.09,0.07)
      node[pos=0, above left=-2pt, align=center] {critical set:\\ $\vect e_R=0$, $\Psi=2$};
  \end{axis}
\end{tikzpicture}

\caption{The attitude error for a rotation of $R_d\transp R$ by the angle
$\theta$ about any unit axis $\vect u$, for which $\tr(R_d\transp R)=1+2\cos\theta$
gives $\Psi=1-\cos\theta$ and $\vect e_R=\sin\theta\,\vect u$. The configuration
error grows monotonically to its maximum $2$ at $\theta=\pi$, but the restoring
term, proportional to $\vect e_R$, peaks at $\pi/2$ and decays to zero at $\pi$
(shaded). Every rotation by $\pi$ is therefore a critical point of $\Psi$ at
which the proportional feedback vanishes, which is why \cref{eq:roa} excludes
$\Psi=2$ and why no continuous feedback of this kind is global on $\SOthree$.}
\label{fig:psi}
\end{figure}

The engine implements
\begin{equation}
  \vect\tau = J\big(-K_R\vect e_R - K_\omega\vect e_\omega\big)
              + \vect\omega\times J\vect\omega
              - J\big(\hatmap{\vect\omega}R\transp R_d\vect\omega_d
                      - R\transp R_d\dt{\vect\omega}{}_d^{\psi}\big),
  \qquad
  \vect e_\omega = \vect\omega - R\transp R_d\vect\omega_d ,
  \label{eq:leeused}
\end{equation}
with the desired angular velocity taken from differential
flatness~\cite{mellinger2011}. With $\vect j_d$ the reference jerk and $T_c$ the
commanded thrust,
\begin{equation}
  \vect h_\omega = \frac{m}{T_c}\Big(\vect j_d
     - (\vect b_3^d\cdot\vect j_d)\,\vect b_3^d\Big),
  \quad
  \vect\omega_d = \begin{bmatrix}
     -\vect h_\omega\cdot\vect b_2^d\\ \vect h_\omega\cdot\vect b_1^d\\
     \dt\psi_d\,(\vect b_3^d\cdot\ez) \end{bmatrix},
  \quad
  \dt{\vect\omega}{}_d^{\psi} = \begin{bmatrix} 0\\0\\
     \ddt\psi_d\,(\vect b_3^d\cdot\ez)\end{bmatrix} .
  \label{eq:omegad}
\end{equation}
\Cref{eq:leeused} differs from the law of \cref{thm:lee} in three ways. The
gains $JK_R$ and $JK_\omega$ are matrices, which coincide with scalar gains only
when $J$ is a multiple of the identity. The tilt part of $\dt{\vect\omega}_d$,
which needs the reference snap, is omitted. And the plant has motor lag and
rotor saturation, which the theorem does not. \Cref{thm:lee} therefore
motivates \cref{eq:leeused} and does not certify it: the behaviour of the
implemented loop is established by simulation for the trajectories reported,
not by proof. Normalising by $J$ makes $K_R$ carry units of
\si{\per\second\squared} and gives it a direct reading: the angular
acceleration commanded per radian of attitude error.

\begin{remark}[Why the rate feedforward matters]
An earlier version of the engine set $\vect\omega_d=0$. On a curved path the
thrust direction rotates at the turn rate of the path (\cref{fig:thrustrot}),
and a loop with no rate
feedforward follows it with a standing attitude error of about
$K_\omega\|\vect\omega_d\|/K_R$. On the turnaround mission of
\cref{sec:results} that version reached a peak tracking error of
\SI{0.27}{\metre} and a rotor disk passed \SI{0.34}{\metre} from the user's
eye. With \cref{eq:omegad} the same mission gives \SI{\rTurnTrack}{\milli\metre}
and \SI{\rTurnClear}{\metre}.
\end{remark}

\begin{figure}[tbp]
\centering
% Why a curved path needs angular-rate feedforward. Left: top view of a turn
% of radius rho at speed V; the horizontal part of b3^d points at the centre
% and turns with the path at Omega = V/rho. Right: b3^d sweeps a cone of
% half-angle theta_t = atan(V^2/(rho g)) about e3, so its derivative has
% magnitude Omega sin(theta_t); omega_d supplies that rate.
\begin{tikzpicture}[font=\footnotesize]
  % ---- left: top view ----
  \begin{scope}
    \def\rr{1.9}
    \draw[construct] (0,0) circle (\rr);
    \fill[cG] (0,0) circle (1.2pt);
    \foreach \a in {90, 30, -30, -90} {
      \coordinate (P) at (\a:\rr);
      \draw[cA, line width=0.6pt] ($(P)+(\a+45:0.16)$) -- ($(P)+(\a+225:0.16)$)
           ($(P)+(\a-45:0.16)$) -- ($(P)+(\a+135:0.16)$);
      \fill[cA] (P) circle (1.3pt);
      \draw[vec, cA] (P) -- ($(P)!0.36!(0,0)$);
      \draw[vec, cB] (P) -- ($(P)+(\a-90:0.62)$);
    }
    \node[cB, anchor=north west] at ($(-90:\rr)+(-180:0.62)+(0.0,-0.02)$) {$\vect v$};
    \node[cA, align=left, anchor=west] at ($(90:\rr)!0.4!(0,0)+(0.1,0.05)$)
         {horizontal part\\ of $\vect b_3^d$};
    % the motion is clockwise (velocity at angle a - 90 degrees)
    \draw[-{Stealth[length=2mm]}, cE, line width=0.8pt]
         (260:0.8) arc (260:190:0.8);
    \node[cE, anchor=east] at (228:0.85) {$\Omega=V/\rho$};
    \draw[dim] (0,0) -- (145:\rr) node[pos=0.6, above, sloped] {$\rho$};
    \node[anchor=north] at (0,-\rr-0.4) {(a) top view of a turn};
  \end{scope}
  % ---- right: the cone swept by b3^d ----
  \begin{scope}[xshift=6.6cm, yshift=-1.7cm]
    \def\L{3.2}
    \def\th{26}
    \pgfmathsetmacro{\rc}{\L*sin(\th)}
    \pgfmathsetmacro{\hc}{\L*cos(\th)}
    \coordinate (O) at (0,0);
    \coordinate (Z) at (0,1);
    \coordinate (S) at (\rc,\hc);
    \draw[vec, cG] (O) -- (0,\L+0.5) node[above] {$\ez$};
    \draw[construct] (0,\hc) ellipse ({\rc} and {0.3*\rc});
    \draw[construct] (O) -- (-\rc,\hc);
    \draw[construct] (O) -- (S);
    \coordinate (B1) at ({\rc*cos(-110)},{\hc+0.3*\rc*sin(-110)});
    \coordinate (B2) at ({\rc*cos(-55)},{\hc+0.3*\rc*sin(-55)});
    \draw[vec, cA!50] (O) -- (B1) node[pos=0.62, left] {$\vect b_3^d(t)$};
    \draw[vec, cA] (O) -- (B2) node[pos=1, below right=-1pt] {$\vect b_3^d(t+\delta t)$};
    \draw[-{Stealth[length=2mm]}, cE, line width=0.8pt, domain=-120:-60, samples=24]
         plot ({1.15*\rc*cos(\x)},{\hc+0.345*\rc*sin(\x)});
    \node[cE, anchor=west, align=left] at ({1.15*\rc},{\hc+0.35})
         {$\|\dt{\vect b}{}_3^d\| = \Omega\sin\theta_t$,\\
          supplied by $\vect\omega_d$ of \labelcref{eq:omegad}};
    \pic[draw, cG, angle radius=1.2cm, "$\theta_t$", angle eccentricity=1.3]
         {angle=S--O--Z};
    \node[anchor=north, align=center] at (0.6,-0.3)
         {(b) $\vect b_3^d$ sweeps a cone about $\ez$,\\ $\tan\theta_t = V^2/(\rho g)$};
  \end{scope}
\end{tikzpicture}

\caption{Why a curved path needs angular-rate feedforward. (a) On a turn of
radius $\rho$ flown at speed $V$ the commanded force has a horizontal part
pointing at the centre of the turn, so the horizontal part of $\vect b_3^d$
turns with the path at $\Omega=V/\rho$. (b) $\vect b_3^d$ therefore sweeps a cone
of half-angle $\theta_t$ about $\ez$, with $\tan\theta_t=V^2/(\rho g)$, and changes
at the rate $\Omega\sin\theta_t$ even in a perfectly steady turn. With
$\vect\omega_d=0$ the attitude loop sees that rate as a persistent error and
settles at a lag of about $K_\omega\|\vect\omega_d\|/K_R$; \cref{eq:omegad}
supplies the rate instead.}
\label{fig:thrustrot}
\end{figure}

\subsection{Control authority and gain selection}

\Cref{thm:lee} says nothing about saturation, and saturation is what actually
limits this vehicle.

\begin{definition}[Torque and angular acceleration authority]
Holding total thrust at $mg$ and the other two moments at zero, the authority
about body axis $j$ is the smaller of the two one-sided maxima
\begin{equation}
  \tau_j^{\max} = \min\{\tau_j^{+},\tau_j^{-}\},
  \qquad
  \tau_j^{\pm} = \max\Big\{\pm\tau :
     \mat M\vect u = (mg,\ \tau\vect e_j),\
     0\le u_i\le\Omega_{\max}^2\Big\} ,
  \label{eq:authority}
\end{equation}
a linear program in $\vect u=(\Omega_1^2,\dots,\Omega_N^2)$ that the engine
solves for every layout, with $\Omega_{\max}^2 = \lambda\,\Omega_h^2$ and
$\Omega_h^2 = mg/(Nk_T)$. For the symmetric quadrotor with alternating spins
it has the closed form (\cref{app:proofs})
\begin{equation}
  \tau_j^{\max} = \Delta\sum_{i=1}^{4}\big|\mat M_{j+1,i}\big|,
  \qquad
  \Delta = \min\{\Omega_h^2,\ \Omega_{\max}^2-\Omega_h^2\} ,
  \label{eq:authquad}
\end{equation}
the differential being bounded by whichever is smaller, the room to speed two
rotors up or the room to slow the other two down. The angular acceleration
authority is $\alpha_j^{\max} = (J^{-1})_{jj}\,\tau_j^{\max}$, which is
$\tau_j^{\max}/J_{jj}$ only when $J$ is diagonal.
\end{definition}

For the design of \cref{sec:results}, \cref{eq:authority} gives
$\tau_{\mathrm{roll}}^{\max}=\tau_{\mathrm{pitch}}^{\max}=\SI{\rTauRoll}{\newton\metre}$
and $\tau_{\mathrm{yaw}}^{\max}=\SI{\rTauYaw}{\newton\metre}$, i.e.\
$\alpha^{\max}_{\mathrm{roll}/\mathrm{pitch}/\mathrm{yaw}} =
\rAlphaRoll/\rAlphaPitch/\rAlphaYaw\ \si{\radian\per\second\squared}$. An earlier
version of \cref{eq:authority} divided the sum by $N$, on the reasoning that
the thrust constraint ``shares'' the differential among the rotors; it does
not, the constraint is met by pairing increments with decrements, and the
error was a factor of four. It was found by independent review and is
recorded in \cref{sec:verification}. \Cref{fig:attainable} draws the whole
attainable set, of which \cref{eq:authority} takes the axis intercepts.

\begin{figure}[tbp]
\centering
% Attainable torque sets at thrust mg and zero yaw moment, computed from the
% engine by one linear program per direction (the support function).
\providecommand{\fdQuadRoll}{3.66}
\providecommand{\fdQuadPitch}{3.66}
\providecommand{\fdQuadOld}{0.91}
\providecommand{\fdHexRoll}{6.13}
\providecommand{\fdHexPitch}{7.07}
\providecommand{\fdHexGross}{2.565}
\providecommand{\fdQuadGross}{1.724}
\providecommand{\fdKeepRadius}{1.180}
\providecommand{\fdKeepSlice}{1.174}
\providecommand{\fdEnvRadius}{1.080}
\providecommand{\fdTurnT}{5.65}
\providecommand{\fdTurnHeadX}{-1.440}
\providecommand{\fdTurnHeadY}{0.000}
\providecommand{\fdTurnPeriod}{7.54}
\providecommand{\fdStandoff}{1.20}
\providecommand{\fdRootLo}{4.5678}
\providecommand{\fdRootHi}{4.6173}
\providecommand{\fdPeakM}{4.5925}
\providecommand{\fdPeakF}{0.03}
\providecommand{\fdGridLo}{4.460}
\providecommand{\fdGridHi}{6.021}
\providecommand{\fdGridLoF}{-0.001}
\providecommand{\fdGridHiF}{-0.083}
\providecommand{\fdIslandPct}{1.1}
\providecommand{\fdPlainRate}{0.263}
\providecommand{\fdConvWindow}{8}
\providecommand{\fdHlMotorRe}{-0.0889}
\providecommand{\fdHlMotorIm}{0.0000}
\providecommand{\fdHlAttRe}{-0.0240}
\providecommand{\fdHlAttIm}{0.0000}
\providecommand{\fdHlPosRe}{-0.0082}
\providecommand{\fdHlPosIm}{0.0004}
\providecommand{\fdHlGovRe}{-0.0160}
\providecommand{\fdHlGovIm}{0.0000}
\providecommand{\fdWpos}{2.04}
\providecommand{\fdWgov}{4.00}
\providecommand{\fdWatt}{6.43}
\providecommand{\fdWmot}{22.2}
\providecommand{\fdWarm}{206}
\providecommand{\fdWnyq}{785}
\providecommand{\fdRatioPA}{3.2}
\providecommand{\fdRatioAM}{3.5}
\providecommand{\fdKR}{41.4}
\providecommand{\fdKp}{4.17}
\providecommand{\fdTauMms}{45}
\providecommand{\fdH}{4}
\providecommand{\fdMotorQ}{0.289}
\providecommand{\fdMotorM}{0.0609}

\begin{tikzpicture}
  \begin{axis}[paperplot, width=0.62\linewidth, height=0.62\linewidth,
      axis equal image,
      xmin=-8, xmax=8, ymin=-8, ymax=8,
      xlabel={$\tau_x$ (roll) [\si{\newton\metre}]},
      ylabel={$\tau_y$ (pitch) [\si{\newton\metre}]},
      legend style={at={(0.5,1.02)}, anchor=south, legend columns=2},
    ]
    \addplot[cA, fill=cA!14] table[x=tx, y=ty] {results/fig/d_attain_quad.dat};
    \addlegendentry{quadrotor, the design}
    \addplot[cB, dashed, line width=1.1pt] table[x=tx, y=ty] {results/fig/d_attain_hex.dat};
    \addlegendentry{hexacopter, same payload and rotors}
    \addplot[only marks, mark=square*, mark size=1.9pt, cB, forget plot]
      coordinates {(\fdHexRoll,0) (-\fdHexRoll,0) (0,\fdHexPitch) (0,-\fdHexPitch)};
    % authority of the quadrotor on each axis
    \addplot[only marks, mark=*, mark size=2.2pt, cA, forget plot]
      coordinates {(\fdQuadRoll,0) (-\fdQuadRoll,0) (0,\fdQuadPitch) (0,-\fdQuadPitch)};
    \node[lbl, cA, anchor=south west] at (axis cs:\fdQuadRoll,0.2)
      {$\tau^{\max}_{\mathrm{roll}}$};
    \node[lbl, cA, anchor=south west] at (axis cs:0.2,\fdQuadPitch)
      {$\tau^{\max}_{\mathrm{pitch}}$};
    % the earlier estimate with the factor 1/N
    \addplot[only marks, mark=x, mark size=3.2pt, cD, line width=1pt, forget plot]
      coordinates {(\fdQuadOld,0) (-\fdQuadOld,0) (0,\fdQuadOld) (0,-\fdQuadOld)};
    \node[lbl, cD, anchor=north west, align=left] at (axis cs:0.15,-\fdQuadOld)
      {earlier estimate\\ with $1/N$: \SI{\fdQuadOld}{\newton\metre}};
    \draw[construct] (axis cs:-8,0) -- (axis cs:8,0);
    \draw[construct] (axis cs:0,-8) -- (axis cs:0,8);
  \end{axis}
\end{tikzpicture}

\caption{Attainable torque sets in the roll and pitch plane with thrust held
at $mg$ and the yaw moment at zero, computed from the engine through their
support function: one linear program of the form \cref{eq:authority} for each
of \num{240} directions. The quadrotor's set is a square turned by \ang{45}
whose axis intercepts are
$\tau^{\max}_{\mathrm{roll}}=\tau^{\max}_{\mathrm{pitch}}=\SI{\fdQuadRoll}{\newton\metre}$,
the balanced differential of \cref{eq:authquad}; the crosses mark the earlier
estimate with the factor $1/N$, a quarter of it. The hexacopter of
\cref{tab:hex}, re-optimised under the same limits (gross
\SI{\fdHexGross}{\kilo\gram}), has a hexagonal set that reaches further in
pitch than in roll: its roll
differential also produces a yaw moment, which the zero-yaw constraint must
cancel, while its pitch differential does not. Its authority in the sense of
\cref{eq:authority} is the smaller intercept, roll.}
\label{fig:attainable}
\end{figure}

\begin{model}[Gain heuristic]
\label{mod:gains}
With $e_{\max}$ the largest attitude error the loop is expected to correct and
$\zeta$ a damping ratio, the engine sets
\begin{equation}
  K_R = \min\left\{\frac{\alpha^{\max}}{2\,e_{\max}},\
                   \Big(\frac{0.3}{\tau_m}\Big)^{2}\right\},
  \qquad
  \alpha^{\max} \coloneqq \min\{\alpha_{\mathrm{roll}}^{\max},
                                \alpha_{\mathrm{pitch}}^{\max}\},
  \qquad
  K_\omega = 2\zeta\sqrt{K_R} ,
  \label{eq:kr}
\end{equation}
and, one level up,
$K_p = \min\{a_{\max}/e_{p,\max},\ K_R/9\}$, $K_d = 2\zeta\sqrt{K_p}$, with
$a_{\max}$ the acceleration cap of the reference governor. The second term
keeps the position bandwidth $\sqrt{K_p}$ at least a factor three below the
attitude bandwidth $\sqrt{K_R}$, which is the time-scale separation the
cascade of \cref{eq:chain} rests on; a machine whose attitude authority has
been cut, by a display on a boom, is destabilised by position gains that were
fine without the boom.
\end{model}

The first bound is the linearised saturation argument: with $\vect e_\omega$
small the commanded angular acceleration is $K_R\|\vect e_R\|\le K_R e_{\max}$,
and half the single-axis authority is reserved for the rate term, for
simultaneous demands on the other axes and for yaw. The second bound keeps the
attitude bandwidth $\sqrt{K_R}$ below $0.3/\tau_m$, so that the first-order
motor lag of \cref{eq:motorlag} contributes less than about \ang{17} of phase
at crossover. $K_\omega$ follows from the linearised closed loop
$\ddt e + K_\omega\dt e + K_R e = 0$. None of this is a saturation certificate:
it neglects the rate feedback term, simultaneous torque demands, the yaw axis
and the actuator dynamics, and \cref{sec:results} verifies by simulation that
the resulting loop stays out of saturation on the missions flown. For the
design the two bounds are \SI{\rKRauth}{\per\second\squared} and
\SI{\rKRlag}{\per\second\squared}, the lag bound binds, and the value used is
$K_R=\SI{\rKR}{\per\second\squared}$, $K_\omega=\SI{\rKw}{\per\second}$.

\begin{figure}[htbp]
\centering
% The frequency ladder of the cascade: each loop is a factor of about three
% or more slower than the one it relies on, and the sampling rate is two
% orders of magnitude above the fastest loop.
\providecommand{\fdQuadRoll}{3.66}
\providecommand{\fdQuadPitch}{3.66}
\providecommand{\fdQuadOld}{0.91}
\providecommand{\fdHexRoll}{6.13}
\providecommand{\fdHexPitch}{7.07}
\providecommand{\fdHexGross}{2.565}
\providecommand{\fdQuadGross}{1.724}
\providecommand{\fdKeepRadius}{1.180}
\providecommand{\fdKeepSlice}{1.174}
\providecommand{\fdEnvRadius}{1.080}
\providecommand{\fdTurnT}{5.65}
\providecommand{\fdTurnHeadX}{-1.440}
\providecommand{\fdTurnHeadY}{0.000}
\providecommand{\fdTurnPeriod}{7.54}
\providecommand{\fdStandoff}{1.20}
\providecommand{\fdRootLo}{4.5678}
\providecommand{\fdRootHi}{4.6173}
\providecommand{\fdPeakM}{4.5925}
\providecommand{\fdPeakF}{0.03}
\providecommand{\fdGridLo}{4.460}
\providecommand{\fdGridHi}{6.021}
\providecommand{\fdGridLoF}{-0.001}
\providecommand{\fdGridHiF}{-0.083}
\providecommand{\fdIslandPct}{1.1}
\providecommand{\fdPlainRate}{0.263}
\providecommand{\fdConvWindow}{8}
\providecommand{\fdHlMotorRe}{-0.0889}
\providecommand{\fdHlMotorIm}{0.0000}
\providecommand{\fdHlAttRe}{-0.0240}
\providecommand{\fdHlAttIm}{0.0000}
\providecommand{\fdHlPosRe}{-0.0082}
\providecommand{\fdHlPosIm}{0.0004}
\providecommand{\fdHlGovRe}{-0.0160}
\providecommand{\fdHlGovIm}{0.0000}
\providecommand{\fdWpos}{2.04}
\providecommand{\fdWgov}{4.00}
\providecommand{\fdWatt}{6.43}
\providecommand{\fdWmot}{22.2}
\providecommand{\fdWarm}{206}
\providecommand{\fdWnyq}{785}
\providecommand{\fdRatioPA}{3.2}
\providecommand{\fdRatioAM}{3.5}
\providecommand{\fdKR}{41.4}
\providecommand{\fdKp}{4.17}
\providecommand{\fdTauMms}{45}
\providecommand{\fdH}{4}
\providecommand{\fdMotorQ}{0.289}
\providecommand{\fdMotorM}{0.0609}

\begin{tikzpicture}
  \begin{axis}[
      width=\linewidth, height=4.6cm, xmode=log,
      xmin=1, xmax=2000, ymin=-1.9, ymax=2.2,
      axis y line=none, axis x line*=bottom,
      xlabel={angular frequency [\si{\radian\per\second}]},
      tick label style={font=\footnotesize}, label style={font=\small},
      clip=false]
    \draw[cG] (axis cs:1,0) -- (axis cs:2000,0);
    \addplot[only marks, mark=*, mark size=2.4pt, cA] coordinates
      {(\fdWpos,0) (\fdWatt,0)};
    \addplot[only marks, mark=square*, mark size=2.2pt, cB] coordinates
      {(\fdWgov,0)};
    \addplot[only marks, mark=triangle*, mark size=2.8pt, cE] coordinates
      {(\fdWmot,0) (\fdWarm,0)};
    \addplot[only marks, mark=diamond*, mark size=2.8pt, cG] coordinates
      {(\fdWnyq,0)};
    \node[lbl, anchor=north, align=center] at (axis cs:\fdWpos,-0.3)
      {position\\ $\sqrt{K_p}=\fdWpos$};
    \node[lbl, anchor=south, align=center, cB] at (axis cs:\fdWgov,0.3)
      {governor\\ $\sqrt{k_p}=4$};
    \node[lbl, anchor=north, align=center] at (axis cs:\fdWatt,-0.3)
      {attitude\\ $\sqrt{K_R}=\fdWatt$};
    \node[lbl, anchor=south, align=center, cE] at (axis cs:\fdWmot,0.3)
      {motor\\ $1/\tau_m=\fdWmot$};
    \node[lbl, anchor=north, align=center, cE] at (axis cs:\fdWarm,-0.3)
      {first arm bending\\ mode, $2\pi f_1=\fdWarm$};
    \node[lbl, anchor=south, align=center, cG] at (axis cs:\fdWnyq,0.3)
      {sampling\\ $\pi/h=\fdWnyq$};
    \draw[dim] (axis cs:\fdWpos,1.55) -- (axis cs:\fdWatt,1.55)
      node[midway, above, font=\scriptsize] {$\times\fdRatioPA$};
    \draw[dim] (axis cs:\fdWatt,1.55) -- (axis cs:\fdWmot,1.55)
      node[midway, above, font=\scriptsize] {$\times\fdRatioAM$};
    \draw[construct] (axis cs:\fdWpos,0.1) -- (axis cs:\fdWpos,1.55);
    \draw[construct] (axis cs:\fdWatt,0.1) -- (axis cs:\fdWatt,1.55);
    \draw[construct] (axis cs:\fdWmot,1.2) -- (axis cs:\fdWmot,1.55);
  \end{axis}
\end{tikzpicture}

\caption{The frequency ladder of the design's cascade on a logarithmic axis.
Each loop is at least a factor of three slower than the one it relies on:
position at $\sqrt{K_p}$, attitude at $\sqrt{K_R}$, the motors at $1/\tau_m$.
The governor's filter at $\sqrt{k_p}$ shapes the reference and closes no loop
on the vehicle. The first bending mode of the arms and the sampling rate of
the controller lie one and two orders of magnitude above the fastest loop,
which is the condition under which the airframe may be treated as rigid and
the sampled controller as continuous.}
\label{fig:bandwidth}
\end{figure}

\begin{remark}[Yaw is the weak axis]
By \cref{eq:alloc} the yaw column involves $k_Q$ rather than $k_T L$, and
$k_Q/k_T \sim \Cdz$-scale quantities. A rotor optimised for quietness by
\cref{obs:sixthpower} turns slowly and has small reaction torque, so
$\alpha_{\mathrm{yaw}}^{\max}$ is an order of magnitude below the roll and
pitch values: \SI{\rAlphaYaw}{\radian\per\second\squared} against
\SI{\rAlphaRoll}{\radian\per\second\squared} for the design. A quiet hovering
display therefore cannot turn to face its user quickly. The engine responds by
governing the heading reference as it governs position, with a rate cap of
\SI{1.5}{\radian\per\second} and an acceleration cap of half the yaw
authority; the architectural response, a yaw axis on the gimbal where a
\SI{30}{\gram} motor solves what the airframe cannot, is not modelled.
\end{remark}

\subsection{Control allocation under saturation}

Given a desired wrench $\vect\chi_d = (T_c,\vect\tau)$ the allocation problem is
\begin{equation}
  \text{find } \vect u\in[0,\Omega_{\max}^2]^N
  \quad\text{with}\quad \mat M\vect u = \vect\chi_d .
  \label{eq:allocprob}
\end{equation}
When $\vect\chi_d$ lies outside $\mat M[0,\Omega_{\max}^2]^N$ no exact solution
exists and a priority must be chosen. We solve the unconstrained least-norm
problem $\vect u = \mat M^{+}\vect\chi_d$ and, on violation, iterate
\begin{equation}
  \vect u \leftarrow \Pi_{[0,\Omega_{\max}^2]}(\vect u),
  \qquad
  \vect u \leftarrow \vect u + \mat M^{+}
    \big(\vect\chi_d - \mat M\vect u\big)\odot \vect\nu ,
  \label{eq:allociter}
\end{equation}
with weights $\vect\nu$ that de-prioritise the thrust channel. The
engineering content is the priority: \emph{attitude is preserved and thrust is
sacrificed}, because losing attitude loses the vehicle whereas losing thrust
loses altitude, which is recoverable.

\subsection{The reference governor}

Tracking a person's face rigidly is the wrong specification, and this is a
statement about kinematics rather than taste.

\begin{proposition}[The cost of a rigid standoff during a turn]
\label{prop:turn}
Let the user rotate their heading by $\pi$ in time $\Delta t$ while the display
is required to remain at standoff $d$ ahead of them. Then the display must
traverse a semicircle of radius $d$, and its mean speed and peak centripetal
acceleration satisfy
\begin{equation}
  \bar v = \frac{\pi d}{\Delta t},
  \qquad
  a_c = \frac{\bar v^2}{d} = \frac{\pi^2 d}{\Delta t^2} .
  \label{eq:turncost}
\end{equation}
\end{proposition}

For $d=\SI{0.7}{\metre}$ and $\Delta t = \SI{1}{\second}$ this is
\SI{2.2}{\metre\per\second} and \SI{6.9}{\metre\per\second\squared}, and the
simulated demand including the translation is over
\SI{10}{\metre\per\second\squared}. The engine confirms that increasing the
thrust margin makes the outcome \emph{worse}, because the controller then tilts
harder and the gimbal cannot keep up.

\begin{model}[Rate-limited reference with a moving keep-out]
\label{mod:governor}
Let $(\vect r_{\mathrm{raw}},\vect v_{\mathrm{raw}},\vect a_{\mathrm{raw}},
\vect j_{\mathrm{raw}})$ be the ideal display trajectory and its derivatives,
and let $\vect c(t)$ be the user's head with velocity $\dt{\vect c}$ and
acceleration $\ddt{\vect c}$. The governed reference is the state
$(\vect r_g,\vect v_g)$ of a second-order filter with acceleration feedforward,
\begin{equation}
  \vect a^{\star} = \vect a_{\mathrm{raw}}
    + k_p(\vect r_{\mathrm{raw}}-\vect r_g)
    + k_d(\vect v_{\mathrm{raw}}-\vect v_g) ,
  \label{eq:gov1}
\end{equation}
whose acceleration is projected, by alternating projections, onto the
intersection of three convex sets: the acceleration ball
$\|\vect a\|\le a_{\max}$, the next-step speed ball
$\|\vect v_g+\vect a\,\dd t\|\le v_{\max}$, and the keep-out half-space
$\vect n\transp\vect a\ge\ell$. With
$\varrho = d_{\mathrm{safe}}+\mathcal{R}+\varepsilon$ the keep-out radius,
$\vect n = (\vect r_g-\vect c)/\|\vect r_g-\vect c\|$,
$h = \|\vect r_g-\vect c\|-\varrho$, $\dt h = \vect n\transp(\vect v_g-\dt{\vect c})$
and $\|\vect v_t\|^2 = \|\vect v_g-\dt{\vect c}\|^2-\dt h^2$,
\begin{equation}
  \ell = \max\left\{
    \vect n\transp\ddt{\vect c} - \frac{\|\vect v_t\|^2}{\|\vect r_g-\vect c\|}
      - 2\gamma\,\dt h - \gamma^2 h ,\;
    \vect n\transp\ddt{\vect c} - \frac{2\,(h + \dt h\,\dd t)}{\dd t^{2}}
  \right\},
  \qquad \gamma = 4 .
  \label{eq:barrier}
\end{equation}
The state then advances by
$\vect r_g\leftarrow\vect r_g+\vect v_g\,\dd t+\tfrac12\vect a\,\dd t^2$,
$\vect v_g\leftarrow\vect v_g+\vect a\,\dd t$. There is no position projection.
If the three sets have empty intersection the step is flagged infeasible, the
acceleration is scaled back to $a_{\max}$, and the violation is reported. The
heading reference is governed separately by the same filter in one dimension
with a rate cap and an acceleration cap.
\end{model}

\begin{proposition}[What the keep-out guarantees]
\label{prop:governor}
\begin{enumerate}[label=(\roman*),nosep]
  \item The caps alone do not keep the reference out of the user. On a curved
        raw path the acceleration cap makes the governed reference cut the
        corner, and since the raw path lies at distance $d$ from the user by
        construction, the governed reference can approach the user
        arbitrarily closely as $a_{\max}$ decreases.
  \item In continuous time, and while the first term of \cref{eq:barrier} is
        the active constraint, the set
        $\{h\ge0,\ \dt h+\gamma h\ge0\}$ is forward invariant: a reference
        that starts outside the keep-out sphere with an admissible closing
        rate never enters it, whatever the user does.
  \item The second term of \cref{eq:barrier} is the one-step condition
        $h_{k+1}\ge0$ under constant $\ddt{\vect c}$ over the step, linearised
        in the tangent plane. It is what the discrete update enforces.
  \item The vehicle, not the reference, is what must stay clear. If
        $\sup_t\|\vect r(t)-\vect r_g(t)\|\le\varepsilon$ on a mission, then on
        that mission the vehicle's centre stays at least
        $d_{\mathrm{safe}}+\mathcal{R}$ from the head, and hence every blade tip
        at least $d_{\mathrm{safe}}$. The hypothesis is checked per run and
        reported; it is not proved.
\end{enumerate}
\end{proposition}

\begin{proof}
(i) In the limit of small $a_{\max}$ the filter output is a convex combination
of past raw positions, and a chord of a circular arc lies inside the arc, so
the reference lies inside the sphere of radius $d$ about the user.
(ii) Differentiating $h$ twice along the governed motion, with
$\dt{\vect n} = \vect v_t/\|\vect r_g-\vect c\|$ in the notation above,
\[
  \ddt h = \vect n\transp(\vect a-\ddt{\vect c})
           + \frac{\|\vect v_t\|^2}{\|\vect r_g-\vect c\|} ,
\]
so the constraint $\vect n\transp\vect a\ge\ell_1$ with $\ell_1$ the first
term of \cref{eq:barrier} is exactly $\ddt h\ge-2\gamma\dt h-\gamma^2 h$. Put
$\psi=\dt h+\gamma h$; then $\dt\psi = \ddt h+\gamma\dt h\ge-\gamma\psi$, so
$\psi(t)\ge\psi(0)e^{-\gamma t}\ge0$, and then
$\dt h\ge-\gamma h$ gives $h(t)\ge h(0)e^{-\gamma t}\ge0$. This is the
second-order barrier construction of Ames et al.~\cite{ames2017}, with both
poles at $-\gamma$ to match the filter's own $k_p=\gamma^2$, $k_d=2\gamma$.
(iii) Expand $\|\vect r_g+\vect v_g\dd t+\tfrac12\vect a\,\dd t^2 -
\vect c-\dt{\vect c}\dd t-\tfrac12\ddt{\vect c}\dd t^2\|$ to first order in the
tangent plane about $\vect n$. (iv) Triangle inequality:
$\|\vect r-\vect c\|\ge\|\vect r_g-\vect c\|-\|\vect r-\vect r_g\|
\ge\varrho-\varepsilon = d_{\mathrm{safe}}+\mathcal{R}$, and the reach bound
\cref{eq:reach}.
\end{proof}

\Cref{fig:barrier} draws part (ii) in the phase plane of $h$.

\begin{figure}[tbp]
\centering
% Phase plane of the keep-out barrier h = ||r_g - c|| - rho. The admissible
% set {h >= 0, hdot + gamma h >= 0} (shaded) is forward invariant under the
% extremal dynamics hddot = -2 gamma hdot - gamma^2 h that the constraint
% allows; a start with hdot + gamma h < 0 is not protected.
% Closed form with psi0 = hdot0 + gamma h0:
%   h(t)    = (h0 + psi0 t) exp(-gamma t)
%   hdot(t) = (hdot0 - gamma psi0 t) exp(-gamma t)
\begin{tikzpicture}
  \pgfmathsetmacro{\gam}{4}
  \begin{axis}[paperplot, height=6.0cm,
      xmin=-0.15, xmax=1.0, ymin=-3.4, ymax=0.9,
      xlabel={$h = \|\vect r_g-\vect c\|-\varrho$ [\si{\metre}]},
      ylabel={$\dt h$ [\si{\metre\per\second}]},
      legend pos=south east]
    % admissible set
    \fill[cC!12] (axis cs:0,0) -- (axis cs:0,0.9) -- (axis cs:1.0,0.9)
                 -- (axis cs:1.0,-4.0) -- cycle;
    \draw[cC!70!black, line width=0.9pt] (axis cs:0,0.9) -- (axis cs:0,0)
         -- (axis cs:0.85,-3.4);
    \node[lbl, cC!60!black, anchor=north east] at (axis cs:0.64,-2.62)
         {$\dt h+\gamma h=0$};
    \node[lbl, cC!60!black, align=center] at (axis cs:0.62,0.45)
         {admissible set\\ $\{h\ge0,\ \dt h+\gamma h\ge0\}$};
    \fill[cD!10] (axis cs:-0.15,-3.4) rectangle (axis cs:0,0.9);
    \node[lbl, cD, rotate=90] at (axis cs:-0.075,-1.25) {inside the keep-out sphere};
    % extremal trajectories from admissible starts
    \foreach \hz/\hd in {0.5/-1.5, 0.8/-3.0, 0.2/-0.6, 0.9/0.5} {
      \addplot[cA, domain=0:2.5, samples=90, variable=\t,
               postaction={decorate}, decoration={markings,
               mark=at position 0.25 with {\arrow{Stealth[length=2mm]}}}]
        ({(\hz + (\hd+\gam*\hz)*\t)*exp(-\gam*\t)},
         {(\hd - \gam*(\hd+\gam*\hz)*\t)*exp(-\gam*\t)});
      \addplot[only marks, mark=*, mark size=1.6pt, cA] coordinates {(\hz,\hd)};
    }
    % an inadmissible start: closing too fast for the distance left
    \addplot[cD, dashed, domain=0:1.2, samples=80, variable=\t,
             postaction={decorate}, decoration={markings,
             mark=at position 0.3 with {\arrow{Stealth[length=2mm]}}}]
      ({(0.3 + (-2.5+\gam*0.3)*\t)*exp(-\gam*\t)},
       {(-2.5 - \gam*(-2.5+\gam*0.3)*\t)*exp(-\gam*\t)});
    \fill[cD] (axis cs:0.3,-2.5) circle (1.6pt);
    \node[lbl, cD, anchor=west] at (axis cs:0.31,-2.62) {$\dt h_0+\gamma h_0<0$};
  \end{axis}
\end{tikzpicture}

\caption{The phase plane of the keep-out barrier $h=\|\vect r_g-\vect c\|-\varrho$
with $\gamma=4$. The shaded set $\{h\ge0,\ \dt h+\gamma h\ge0\}$ is forward
invariant under the slowest recession that the first term of
\cref{eq:barrier} permits, $\ddt h=-2\gamma\dt h-\gamma^2h$, whose solutions
$h(t)=(h_0+\psi_0t)\,e^{-\gamma t}$ with $\psi_0=\dt h_0+\gamma h_0$ are drawn
from four admissible starts: each reaches the origin without crossing $h=0$.
A start that closes faster than $\gamma h_0$ (dashed) is not protected and
enters the sphere. This is why the governor maintains the invariant set and
not the half-space $h\ge0$ alone.}
\label{fig:barrier}
\end{figure}

Three things the proposition does not say. The invariance in (ii) is a
statement about the continuous-time reference; the engine enforces (iii) at
update instants, and between them the reference moves under a constant
acceleration, so (ii) is the limit the discrete scheme approaches, not what it
executes. The guarantee lapses whenever the step is infeasible, which is why
that is reported as a violation rather than absorbed. And $\vect c$ is the
true head position: an estimator with latency or dropouts moves the sphere
away from the head, and no estimator is modelled (\cref{as:estimation}).

\Cref{fig:turnaround} shows the governor on the most demanding mission of
\cref{sec:results}, the turnaround.

\begin{figure}[tbp]
\centering
% The turnaround flown by the engine. (a) Top view around one reversal:
% the user's head, the ideal display trajectory (a semicircle about the
% head), the governed reference and the vehicle, with the keep-out circle at
% the reversal. (b) Distances from the head over the whole window.
\providecommand{\fdQuadRoll}{3.66}
\providecommand{\fdQuadPitch}{3.66}
\providecommand{\fdQuadOld}{0.91}
\providecommand{\fdHexRoll}{6.13}
\providecommand{\fdHexPitch}{7.07}
\providecommand{\fdHexGross}{2.565}
\providecommand{\fdQuadGross}{1.724}
\providecommand{\fdKeepRadius}{1.180}
\providecommand{\fdKeepSlice}{1.174}
\providecommand{\fdEnvRadius}{1.080}
\providecommand{\fdTurnT}{5.65}
\providecommand{\fdTurnHeadX}{-1.440}
\providecommand{\fdTurnHeadY}{0.000}
\providecommand{\fdTurnPeriod}{7.54}
\providecommand{\fdStandoff}{1.20}
\providecommand{\fdRootLo}{4.5678}
\providecommand{\fdRootHi}{4.6173}
\providecommand{\fdPeakM}{4.5925}
\providecommand{\fdPeakF}{0.03}
\providecommand{\fdGridLo}{4.460}
\providecommand{\fdGridHi}{6.021}
\providecommand{\fdGridLoF}{-0.001}
\providecommand{\fdGridHiF}{-0.083}
\providecommand{\fdIslandPct}{1.1}
\providecommand{\fdPlainRate}{0.263}
\providecommand{\fdConvWindow}{8}
\providecommand{\fdHlMotorRe}{-0.0889}
\providecommand{\fdHlMotorIm}{0.0000}
\providecommand{\fdHlAttRe}{-0.0240}
\providecommand{\fdHlAttIm}{0.0000}
\providecommand{\fdHlPosRe}{-0.0082}
\providecommand{\fdHlPosIm}{0.0004}
\providecommand{\fdHlGovRe}{-0.0160}
\providecommand{\fdHlGovIm}{0.0000}
\providecommand{\fdWpos}{2.04}
\providecommand{\fdWgov}{4.00}
\providecommand{\fdWatt}{6.43}
\providecommand{\fdWmot}{22.2}
\providecommand{\fdWarm}{206}
\providecommand{\fdWnyq}{785}
\providecommand{\fdRatioPA}{3.2}
\providecommand{\fdRatioAM}{3.5}
\providecommand{\fdKR}{41.4}
\providecommand{\fdKp}{4.17}
\providecommand{\fdTauMms}{45}
\providecommand{\fdH}{4}
\providecommand{\fdMotorQ}{0.289}
\providecommand{\fdMotorM}{0.0609}

\begin{tikzpicture}
  \pgfmathsetmacro{\ta}{\fdTurnT-1.5}
  \pgfmathsetmacro{\tb}{\fdTurnT+2.2}
  \begin{axis}[paperplot, name=top, width=0.47\linewidth,
      axis equal image, xmin=-3.0, xmax=2.3, ymin=-1.35, ymax=1.65,
      xlabel={$x$ [\si{\metre}]}, ylabel={$y$ [\si{\metre}]},
      title={(a) top view, $t\in[\pgfmathprintnumber[fixed,precision=1]{\ta},
             \pgfmathprintnumber[fixed,precision=1]{\tb}]\,\si{\second}$},
      window/.style={restrict expr to domain={\thisrow{t}}{\ta:\tb},
                     unbounded coords=jump}]
    \addplot[cE, line width=2pt, window] table[x=hx, y=hy] {results/fig/d_turn.dat};
    \addplot[cG, dashed, window] table[x=rawx, y=rawy] {results/fig/d_turn.dat};
    \addplot[cB, window] table[x=govx, y=govy] {results/fig/d_turn.dat};
    \addplot[cA, window, postaction={decorate}, decoration={markings,
             mark=between positions 0.12 and 0.95 step 0.2 with {\arrow{Stealth[length=2.2mm]}}}]
      table[x=vehx, y=vehy] {results/fig/d_turn.dat};
    \addplot[cD, dashed, domain=0:360, samples=120, line width=0.7pt]
      ({\fdTurnHeadX + \fdKeepSlice*cos(x)}, {\fdTurnHeadY + \fdKeepSlice*sin(x)});
    \addplot[only marks, mark=*, mark size=2.2pt, cE] coordinates {(\fdTurnHeadX,\fdTurnHeadY)};
    \node[lbl, cE, anchor=north] at (axis cs:\fdTurnHeadX,-0.08) {head at the reversal};
  \end{axis}
  \begin{axis}[paperplot, name=dist, at={(top.east)}, anchor=west, xshift=1.7cm,
      width=0.47\linewidth, height=4.9cm,
      xlabel={$t$ [\si{\second}]}, ylabel={distance from the head [\si{\metre}]},
      title={(b) distances over the window},
      xmin=0, xmax=16, ymin=1.0, ymax=1.35,
      legend to name=turnleg, legend columns=3,
      legend style={/tikz/every even column/.append style={column sep=6pt}}]
    \addlegendimage{cE, line width=2pt}
    \addlegendentry{head}
    \addplot[cG, dashed] table[x=t, y=draw] {results/fig/d_turn.dat};
    \addlegendentry{ideal display trajectory}
    \addplot[cB] table[x=t, y=dgov] {results/fig/d_turn.dat};
    \addlegendentry{governed reference $\vect r_g$}
    \addplot[cA] table[x=t, y=dveh] {results/fig/d_turn.dat};
    \addlegendentry{vehicle $\vect r$}
    \addplot[cD, dashed, domain=0:16] {\fdKeepRadius};
    \addlegendentry{keep-out radius $\varrho=d_{\mathrm{safe}}+\mathcal{R}+\varepsilon$}
    \addplot[cC, dotted, line width=1.3pt, domain=0:16] {\fdEnvRadius};
    \addlegendentry{$d_{\mathrm{safe}}+\mathcal{R}$}
    \draw[construct] (axis cs:\fdTurnT,1.0) -- (axis cs:\fdTurnT,1.35);
  \end{axis}
  \node[anchor=north] at ({$(top.outer west)!0.5!(dist.outer east)$} |- dist.outer south)
       {\pgfplotslegendfromname{turnleg}};
\end{tikzpicture}

\caption{The turnaround flown by the engine with the governor of
\cref{mod:governor} at $v_{\max}=\SI{2.5}{\metre\per\second}$,
$a_{\max}=\SI{2.5}{\metre\per\second\squared}$ and
$\varepsilon=\SI{0.10}{\metre}$. (a) Top view around one reversal of the user.
The ideal display trajectory is a semicircle of radius
$d=\SI{\fdStandoff}{\metre}$ about the head, which \cref{prop:turn} shows is
not flyable; the governed reference yields, swings wide and catches up after
the turn, and the vehicle follows it, the arrows giving the direction of
travel. The dashed circle is the keep-out sphere at the reversal, cut at the
vehicle's altitude. (b) Distances from the head over the whole window, the
vertical line marking the reversal shown in (a). The governed reference does
not enter $\varrho$, as \cref{prop:governor}(ii) and (iii) require; the
vehicle falls below $\varrho$ by at most its tracking error and stays above
$d_{\mathrm{safe}}+\mathcal{R}$, which is \cref{prop:governor}(iv) with the
tracking bound measured on this run. The peaks are the display lagging the
user after each reversal.}
\label{fig:turnaround}
\end{figure}

\begin{remark}[The governor's bandwidth must be high]
The gains $k_p,k_d$ in \cref{eq:gov1} must place the filter's natural frequency
well above the frequency content of the mission; otherwise the governor adds a
steady lag on every curved path, which is a different and worse behaviour than
yielding during a turn. With $k_p=16$, $k_d=8$ the filter is critically damped
at \SI{4}{\radian\per\second}, and with the feedforward $\vect a_{\mathrm{raw}}$
in \cref{eq:gov1} it is exactly transparent whenever no cap binds. Without the
feedforward it lags every acceleration by $\vect a_{\mathrm{raw}}/k_p$, about
\SI{5}{\centi\metre} at \SI{0.8}{\metre\per\second\squared}, which was enough
to push the reference into the keep-out sphere a third of the time on a mission
with no turn in it.
\end{remark}

\begin{definition}[Two error measures]
With a governor in the loop, two distinct errors must be reported:
\begin{equation}
  e_{\mathrm{track}} = \|\vect r_g - \vect r\|
  \quad\text{(control performance)},
  \qquad
  e_{\mathrm{lag}} = \|\vect r_{\mathrm{raw}} - \vect r\|
  \quad\text{(what the user sees)} .
  \label{eq:twoerrors}
\end{equation}
Reporting only $e_{\mathrm{track}}$ would flatter the design, since the governor
is free to make its own reference easy to follow. The peak of
$e_{\mathrm{track}}$ over the whole window, start-up included, is the quantity
that must stay below $\varepsilon$ for \cref{prop:governor}(iv).
\end{definition}

\subsection{Hover linearisation and controllability}

Linearising \cref{eq:ode} about the hover trim by central differences gives
$\dt{\delta x} = \mat A\,\delta x + \mat B\,\delta u$ on the $12+N$ rigid-body
and rotor states, after removing the quaternion scalar component which is not
a free coordinate. The trim is computed, not assumed: with the display on a
boom the centre of mass is off the rotor centroid and hover needs unequal
rotor speeds. Controllability is assessed by the rank of the reachable
subspace, computed by an orthogonal iteration
$Q_{k+1} = \operatorname{range}[\mat B\ \ Q_k\ \ \mat A Q_k]$ on the matrices
scaled by their 2-norms with a singular-value tolerance of $10^{-9}$, which
avoids forming the ill-conditioned powers of the Kalman matrix. Controllability
requires the rank to equal the state dimension; for the design of
\cref{sec:results} it is $16$ of $16$. The eigenvalues cluster at the origin,
reflecting the six open-loop integrators of \cref{eq:chain}, which is
precisely why the cascade rather than a single loop is required.
